\documentclass[11pt]{article}
\usepackage{fullpage}
%\usepackage{setspace}
\usepackage{enumitem}
\usepackage{listings,amssymb}
\usepackage{listings}
\usepackage{graphicx,xcolor}
%\doublespacing
\usepackage[T1]{fontenc}
\usepackage{newtxtext,newtxmath}
\usepackage{hyperref}
%\usepackage[colorlinks=true, urlcolor=blue, linkcolor=blue, citecolor=blue]{hyperref}
\lstset{numbers=left, numberstyle=\tiny}
\title{CS 147: Using the Synthesizer}
\author{}
\date{\today}
\begin{document}
\maketitle
%\vspace{-0.5in}
\section{Overview}
In this class we will do some very simple synthesis of your designs. The primary goal of this exercise is to get a sense for the actual hardware your Verilog is creating. Synthesizing your design will allow us to see:
\begin{itemize}
\item
The actual gates
\item
Area
\item
Delay
\item
The actual longest path in your design :)
\end{itemize}
There is a lot more to synthesis optimizations than what we will cover in this class.
We will be using Synopsys DC Compiler and a 45nm gate library provided by FreePDK. A lot of the details will be abstracted away and you will be using a simple script called \texttt{synth.pl} which will do most of the work for you.

\section{First Time Setup}
In the Git repo for the course, you will find the following folder: \ \texttt{assignments/tools/synth}.  From this folder, copy
the setup file to your home directory on one of the \texttt{coe-ee} machine:
\begin{scriptsize}
\begin{verbatim}
prompt> scp -p assignments/tools/synth/.synopsys.dc_setup <your Okta ID>@coe-ee-cad15.sjsuad.sjsu.edu:
\end{verbatim}
\end{scriptsize}

\textbf{Note:} There are about 50 or so machines in the CAD lab.  Pick a machine that is lightly loaded for your work.

Let \texttt{code} be the folder on the \texttt{coe-ee} machine for doing the work. The Git repo contains a
wrapper script that simplifies the interface to the Synopsis DC compiler which we use for Synthesis.

\begin{scriptsize}
\begin{verbatim}
prompt> scp -p assignments/tools/synth/synth.pl <your Okta ID>@coe-ee-cad15.sjsuad.sjsu.edu:code
\end{verbatim}
\end{scriptsize}


\section{Running the Synthesizer}
On an on-going basis, whenever you need to run the synthesizer, the first step is to copy over the associated files.  For instance, to run the synthesizer
on the Verilog files from HW2.2 (i.e., the ALU), do the following:

\begin{enumerate}
\item Copy the Verilog files
\begin{small}
\begin{verbatim}
prompt> scp -rp assignments/hw02 <your Okta ID>@coe-ee-cad15.sjsuad.sjsu.edu:code
\end{verbatim}
\end{small}
\item Login to the Linux machine and get setup:
\begin{small}
\begin{verbatim}
prompt> ssh <your Okta ID>@coe-ee-cad15.sjsuad.sjsu.edu

coe-ee-cad15> csh

coe-ee-cad15>  source /apps/design_environment.csh

coe-ee-cad15>  cd code/hw02/hw2_2
\end{verbatim}
\end{small}
\item
Make a list of Verilog files that are part of the design. Create a text file (named \texttt{list.txt} in this folder (\texttt{code/hw02/hw2\_2})) 
with one Verilog filename per line.  Additonally,
\begin{itemize}
\item No testbenches should be included in this list
\item No \_hier files which are wrappers should be in this list
\item clkrst.v should NOT be in the list, since we don't want to synthesize the clock generator
\end{itemize}
For example \texttt{list.txt} with the following contents.
\begin{verbatim}
alu.v
4_1mux.v
CLA.v
\end{verbatim}
\item
Check the design.
\begin{small}
\begin{verbatim}
coe-ee-cad15> perl synth.pl --list=list.txt --type=other --cmd=check --top=alu
\end{verbatim}
\end{small}
Look at the output on the screen and \texttt{synth.log}.  If there are no errors, go to the next step.
\item 
Run the synthesis.
\begin{small}
\begin{verbatim}
coe-ee-cad15> perl synth.pl --list=list.txt --type=other --cmd=synth --top=alu
\end{verbatim}
\end{small}
This step usually takes a few minutes.  The output is in \texttt{synth/} folder.  You will find
an area report, cell report, and a timing report.
\item
Optimizing your design.  This is analogous to using the GCC compiler to optimize code generation.
\begin{small}
\begin{verbatim}
coe-ee-cad15> perl synth.pl --list=list.txt --type=other --cmd=synth --top=alu --opt
\end{verbatim}
\end{small}
\end{enumerate}

\section{Interpreting the Output Files}
Please see \href{https://karu.sites.cs.wisc.edu/courses/cs552/spring2021/wiki/index.php/Synthesis/Synthesis#toc11}{this page} about how
to interpret the outputs in the files generated by the DC compiler.

\section{Acknowledgements}
The script \texttt{synth.pl} was provided by Prof. Karu Sankaralingam from Computer Sciences Department of the University of Wisconsin at Madison.
We thank him for his support.

\end{document}